\section{The projective-complement closure of a logic}
\label{sec:p03-delta-closure}

The closure considered here is the Souslin--Kleene closure, defined by
complementary projective classes. Projective witnesses may have larger
universes. The construction therefore retains the witnesses' domains,
not only their additional relation symbols. It uses the guarded
translation interface of \cref{def:p03-w03-p03-048-b74-logic}, with the
first-order and quantifier closures stated there, on the full category
of finite signatures. This is the setting of the closure problem in
\cite[Part~II, \S4, pp.~260--262]{Barwise1974Axioms}, with the
composition qualification already proved in
\cref{thm:p02-translation-composition}.

\begin{hbdefinition}{def:p03-delta-syntax}{Complementary projective indices}
For a finite relational signature \(\rho\), an index consists of two
relativized projective presentations in \(\mathcal L\) whose classes
are complementary among all \(\rho\)-structures. Each presentation
has a finite witness signature and one sentence. Write \(K\) for the
positive class. The generalized quantifier \(Q_K\), applied to a list
of formulas of the arities in \(\rho\), asserts that the structure on
the current universe with exactly those defined relations belongs to
\(K\). Further free variables in those formulas act as parameters.

The syntax of \(\Delta(\mathcal L)\) is finite first-order syntax,
with these additional quantifiers. Each formula uses finitely many
indices, each supplied with its two witness presentations. The assertion
that a proposed pair is complementary is a semantic condition on an
index, not an asserted effective test. Function signatures are handled
by their relational graphs and the usual totality and uniqueness
conditions. Projective classes and complements use the convention of
\cref{def:p03-adm-relativized-pc}, without restricting the ambient
witness to the reduct's universe.
\end{hbdefinition}

\begin{hblemma}{thm:p03-delta-packing}{proved=annex:p3.delta.1}{Packing parameterized witnesses}
A finite formula built with the preceding generalized quantifiers has
both a positive and a negative relativized projective presentation in
\(\mathcal L\). The same holds for formulas with a fixed finite list
of parameters, regarded as named elements of the reduct. More generally,
a finite list of such witness requirements, for every tuple satisfying
a first-order definable activation condition, can be expressed by one
\(\mathcal L\)-sentence in a finite expanded signature.
\end{hblemma}

\begin{hbtheorem}{thm:p03-delta-closure}{proved=annex:p3.delta.2}{The closure theorem}
Up to equivalence on each finite signature:
\begin{enumerate}
\item \(\mathcal L\leq\Delta(\mathcal L)\);
\item \(\mathcal L\) and \(\Delta(\mathcal L)\) have exactly the
      same relativized projective classes;
\item a class is definable in \(\Delta(\mathcal L)\) exactly when
      it and its complement are projective in \(\mathcal L\);
\item \(\Delta(\mathcal L)\) has the Souslin--Kleene property and
      is the least extension of \(\mathcal L\) with that property.
\end{enumerate}
Consequently this closure is monotone and idempotent. These conclusions
do not identify Souslin--Kleene closure with the full Beth property for
possibly partial implicitly defined relations.
\end{hbtheorem}
